\addtocontents{toc}{\protect\pagebreak}

\chapter{Observations aberrantes et influentes}
\label{sec:DonneesInfluentes}

\vspace{-0.1in}
%Tout comme on a sélectionné les variables à inclure dans le modèle de RLM, il peut être nécessaire de sélectionner les observations à inclure dans le modèle.
On doit parfois retirer du modèle une ou plusieurs observations qui ont une influence trop grande sur l'estimation des paramètres. Avant de parler plus en détails de ces données influentes, il faut bien définir deux termes importants.

\begin{defn}\label{def:aber}
Une observation est \textbf{ aberrante} si elle ne cadre pas avec le modèle.
\end{defn}

\vspace{-0.1in}
Ce sera donc une observation qui a un résidu de grande taille. Une donnée aberrante n'est pas problématique en elle-même, mais elle peut indiquer une erreur dans la valeur d'une variable exogène ou de la variable endogène.

\begin{defn}\label{def:infl}
Une donnée \textbf{influente} est une observation qui, \`a elle seule, peut avoir un impact sérieux sur les conclusions d'une analyse de régression.
\end{defn}

\vspace{-0.1 in}
Parmi les exemples \og d'impact sérieux \fg, on peut citer des $\hat{\beta}$ négatifs quand on sait qu'ils devraient \^etre positifs (ou vice-versa) ou encore une variable exogène qui devrait \^etre importante mais qui est considérée non significative. Une ou des observations influentes peut également diminuer le pouvoir prédictif du modèle.

%\newpage
\section{Sources de l'influence}

L'influence d'une observation provient souvent d'une combinaison de
deux sources :
\begin{description}
\item[Source 1:] L'observation a un grand levier, c'est-à-dire qu'elle a une grande capacité \textit{a priori} à influencer les résultats de la régression.
\item[Source 2:] L'observation a une valeur extrême de $Y$ par rapport au modèle.
%Donnée ayant une erreur produisant une valeur extr\^eme de $Y$.
\end{description}
On peut facilement identifier les observations avec ces deux caractéristiques.

\textbf{Levier d'une observation}

Une observation a un grand levier si sa valeur de $\boldx_i$ est loin du centre (centro\"{\i}de) de toutes les observations. La distance entre une observation et le centro\"{\i}de de toutes les autres peut être calculée facilement à partir de la matrice d'incidence $\boX$. Plus précisément, le levier de la $i^e$ observation est le $i^e$ élément diagonal de la matrice chapeau $\boH=\boX(\boX'\boX)^{-1}\boX'$, qui peut s'exprimer comme
$h_{ii}=\boldx_i'(\boX'\boX)^{-1}\boldx_i$.

Qu'est-ce qu'une grande valeur de $h_{ii}$? Il est difficile de
répondre \`a cette question de fa\c{c}on générale. Cependant nous
savons que tr$(\boH)=\sum_{i=1}^n h_{ii}=p'$, o\`u $p'$ est le nombre
de paramètres dans le modèle. Donc un $h_{ii}$ moyen prend une valeur de $p'/n$
et un $h_{ii}$ qui se démarque nettement de cette valeur
signifie une donnée \`a grand levier. \citet{Bel_Kuh_Wel:80} suggèrent de retenir les valeurs de $h_{ii}>2p'/n$.

\textbf{Observation extrême}

Une observation est extrême (ou aberrante) par rapport au modèle si son résidu est grand. On regardera plus précisément le résidu studentisé. On rappelle que celui-ci est donné par

$$T_i={\hat\varepsilon_i\over \hat\sigma_{-i}\sqrt{1-h_{ii}}} \sim\mathcal{T}_{n-p-2},$$
o\`u $$\hat\sigma_{-i}=\sqrt{{(n-p)\hat\sigma^2-\hat\varepsilon_i^2/(1-h_{ii})\over n-p-1}}.$$

Comme on peut le voir, ce résidu est une fonction croissante à la fois en
$h_{ii}$ et en $\hat\varepsilon_i$.

\citet{Bel_Kuh_Wel:80} prônent de prendre les valeurs de $|T_i|>t_{n-p',(1-\alpha^*/2)}$, où $\alpha^*=\alpha/n$ (le seuil corrigé de Bonferroni) comme un bon indicateur d'une valeur aberrante. D'autres statisticiens suggèrent de retenir les valeurs de $|T_i|>2$.

\begin{remarque}
Avoir un grand levier et être une observation extrême sont des indices qu'une observation pourrait avoir une influence démesurée sur les résultats de la régression. Pour tester si cela est effectivement le cas, on doit mesurer directement l'influence de l'observation sur la régression. Il est important de comprendre qu'une donnée avec un grand levier peut ne pas être influente. Il en est de même pour une observation extrême.
\end{remarque}


%Un grand levier n'implique pas nécessairement une influence. Tout diagnostic d'influence doit \^etre basé sur une combinaison de $h_{ii}$ et $R_{i}$.

\section{Mesures d'influence}

Pour mesurer si une donnée est influente ou non, on évalue son impact sur la prévision $\hY$ et sur l'estimation de $\bobeta$. Les différentes statistiques utilisées sont présentées ici.

\subsection{Influence sur la prévision}
La statistique suivante mesure l'influence de la $i^e$ donnée
dans la prévision $\hat{Y}_i$ de $Y_i$ (\textit{Difference in Fit}) :
$$DFFIT_i={\hat{Y}_i-\hat{Y}_{i,-i}\over
\hat\sigma_{-i}\sqrt{h_{ii}}}=T_i\sqrt{h_{ii}\over1-h_{ii}},$$ o\`u
$\hat{Y}_{i,-i}$ désigne la valeur de la prévision de $Y_i$ basée
sur un jeu de données qui ne contient pas la $i^e$ donnée,
c'est-à-dire $\hat{Y}_{i,-i}=\boldx_i'\hat{\bobeta}_{-i}$.

Il faut se rappeler ici que $\var[\hat Y_i]=\sigma^2h_{ii}$. Les différences au numérateur de $DFFIT_i$ sont donc standardisées de sorte qu'elles s'expriment en multiples de l'écart-type de la prévision.

Nous analysons les $DFFIT_i$ de la m\^eme fa\c{c}on que nous nous penchons sur les $T_i$, c'est-à-dire que des grandes valeurs de $|DFFIT_i|$ constituent un signe que la $i^e$ donnée possède une grande influence sur la valeur de la prévision. \citet{Bel_Kuh_Wel:80} indiquent qu'il faut retenir les valeurs de $|DFFIT_i|>2\sqrt{p/n}$ alors que d'autres statisticiens proposent simplement de retenir les valeurs de $|DFFIT_i|>2$.

\subsection{Influence sur les estimations des paramètres}
La statistique suivante mesure l'influence de la $i^e$
observation sur la valeur de l'estimateur de $\beta_j$ (\textit{ Difference in Beta}):
$$DFBETA_{ji}={\hat{\beta}_j-\hat{\beta}_{j,-i}\over
\hat\sigma_{-i}\sqrt{(\boX'\boX)^{-1}_{jj}}}.$$
Il faut se souvenir ici que $\var[\hat\beta_j]=\sigma^2(\boX'\boX)^{-1}_{jj}$. Les différences au numérateur de la statistique sont  à nouveau standardisées.
Encore une fois, les grandes valeurs de $DFBETA_{ji}$ indiquent que la $i^e$ donnée possède une grande influence sur l'estimation de $\beta_j$. \citet{Bel_Kuh_Wel:80} proposent de considérer les valeurs de $|DFBETA_{ji}|>2/\sqrt{n}$ alors que d'autres statisticiens suggèrent simplement de retenir les valeurs de $|DFBETA_{ji}|>2$. \\

Le signe de $DFBETA_{ji}$ est un bon indicateur. En effet, si un certain $\hat\beta_j$ est négatif alors qu'on se serait attendu à une valeur positive, alors un $DFBETA_{ji}$ très négatif indique que le \og coupable \fg pourrait être l'observation $i$.

Comme nous obtenons $p$ valeurs de $DFBETA$
pour chaque observation, cette mesure devient peu pratique lorsque
nous sommes en présence d'un jeu de données ayant plusieurs
variables exogènes. Pour simplifier la recherche de valeurs influentes, Dennis Cook a proposé la statistique suivante, appelée distance
de Cook, qui permet de combiner les $DFBETA$ d'une m\^eme observation:
\begin{eqnarray*}
\text{distance de Cook}_i\;=\;C_i & = & {(\hat{\bobeta}-\hat{\bobeta}_{-i})'(\boX'\boX)(\hat{\bobeta}-\hat{\bobeta}_{-i})
\over p' \hat\sigma^2} \\
& = & \frac{(\hat Y_i-\bo{x}_i'\hat{\bobeta}_{-i})^2}{p'\hat\sigma^2} \\
& = & \frac{\hat\varepsilon_i^2}{(1-h_{ii})^2}\frac{h_{ii}}{\hat\sigma^2 p'} \\
% & = & \frac{\hat\varepsilon_{i,-i}^2}{p}\frac{h_{ii}}{1-h_{ii}}
\end{eqnarray*}
 Une grande valeur de $C_i$ indique une observation
ayant une influence prépondérante sur les $\beta$. \citet{Fox:91} suggère de porter notre attention sur les valeurs de $C_i>\frac{4}{n-p}$. Il s'agit par la suite
d'examiner les $DFBETA$ pour cette observation pour savoir quel(s)
$\hat\beta$ est (sont) affecté(s).


\subsection{Influence sur la performance des estimateurs}
Une donnée peut aussi avoir un impact néfaste sur la variance des estimateurs. Soit $||Var(\hat{\bobeta})|| = (\sigma^2)^p||(\boX'\boX)^{-1}||$, le déterminant de la matrice de variance de $\hat{\bobeta}$, parfois appelé variance généralisée. L'influence de la $i^e$ donnée sur la variance de $\hat{\bobeta}$ peut \^etre calculée comme suit:

\begin{eqnarray*}COVRATIO_i&=&{||(\boX_{-i}'\boX_{-i})^{-1}\hat\sigma_{-i}^2||
\over||(\boX'\boX)^{-1}\hat\sigma^2||}\\
&=&\left({\hat\sigma_{-i}^{2}\over
\hat\sigma^{2}}\right)^p\left({h_{ii}\over1-h_{ii}}\right).\end{eqnarray*}
Comme la seconde ligne l'illustre, une donnée possédant beaucoup de levier
donnera en général un $COVRATIO>1$, et une donnée aberrante un $COVRATIO<1$.

\citet{Bel_Kuh_Wel:80} proposent de porter notre attention sur les valeurs telles que
$|COVRATIO-1|~>~3p/n$, c'est-à-dire sur les observations telles que\\ $COVRATIO~<~1~-~3p/n$ ou $COVRATIO>1+3p/n$.



\section{Que faire avec les données influentes?}
Si une donnée est influente, il faut prendre tous les moyens
possibles pour essayer de comprendre comment cette donnée a été
obtenue et s'assurer qu'il n'y a pas d'erreur qui aurait pu
survenir lors de la mesure, collecte ou saisie des données.

Si l'influence provient d'un grand
levier, alors il serait souhaitable d'aller chercher d'autres
données afin de combler le vide entre le centro\"{\i}de des
données et la donnée influente.

En général, on doit s'abstenir d'éliminer les données
influentes d'un jeu de données. Quelques exceptions:
\begin{itemize}
\item Erreur de mesure ou  de  saisie (corriger la donnée
est préférable \`a son élimination). \item Valeurs
mesurées scientifiquement peu plausibles. \item Donnée très
peu représentative du reste des données et de la population en
général. \end{itemize}

Il est \`a noter que souvent, une transformation de variable peut
grandement réduire l'influence de certaines données. 